\documentclass[10pt,a4paper]{article}
\usepackage[T1]{fontenc}
\usepackage{lmodern,amsmath,amssymb,amsthm,booktabs,graphicx,microtype}
\usepackage[margin=24mm]{geometry}
\usepackage[colorlinks=true,linkcolor=blue,urlcolor=blue,citecolor=blue]{hyperref}
\newtheorem{theorem}{Theorem}
\newtheorem{lemma}{Lemma}
\newtheorem{corollary}{Corollary}
\setlength{\parskip}{1.5pt}
\hypersetup{pdftitle={An explicit formula for Campbell's variable-depth recurrence},
pdfauthor={}}
\title{An explicit formula for Campbell's\texorpdfstring{\\}{ }variable-depth recurrence}
\author{}
\date{29 September 2026}
\begin{document}
\maketitle
\begin{abstract}
We give an explicit formula for the nested recurrence proposed by John M.
Campbell, in which the number of iterations is itself a preceding sequence
term. Its even and odd subsequences consist of alternating linear pieces and
plateaux on powers-of-three scales. The proof identifies every prescribed
inner orbit after four steps: it is then fixed or two-periodic, and the
iteration depth selects the required phase. We also obtain exact dilation
and the lower and upper limits of the ratio of the sequence to its index.
\end{abstract}

\section{The recurrence and its explicit formula}
Campbell's proposed recurrence\footnote{John M. Campbell,
\texttt{MetaMetaV32.pdf}, unpublished note communicated 29 September 2026.} is
\begin{equation}\label{eq:recurrence}
b(1)=1,\qquad T_n(x)=n-b(x),\qquad
b(n)=T_n^{\,b(n-1)}(n-1)\quad(n\ge2).
\end{equation}
Here the exponent denotes iteration of the map $x\mapsto n-b(x)$ with $n$
fixed. The first terms are
\[
1,1,2,3,2,3,4,5,6,7,6,9,6,9,6,9,8,9,10,\ldots.
\]
\begin{theorem}\label{thm:main}
The recurrence defines a unique positive integer sequence. For $n\ge2$,
let $s$ be the unique power of $3$ in the applicable interval below. Then
\begin{equation}\label{eq:formula}
\begin{aligned}
b(n)&=\min(n-s,3s),&& n\text{ even},&2s\le n<6s,\\
b(n)&=\max(2s,n-3s),&& n\text{ odd},&3s\le n<9s.
\end{aligned}
\end{equation}
For every $n\ge2$, the orbit starting at $n-1$ satisfies
\begin{equation}\label{eq:orbit}
T_n^6(n-1)=T_n^4(n-1).
\end{equation}
Its preperiod is at most four and its eventual period is one or two.
The bound four is attained.
\end{theorem}

For example, at scale $s=9$ the even subsequence follows $n-9$ from
$n=18$ through $36$, and then stays at $27$ through $52$. The odd
subsequence stays at $18$ from $27$ through $45$, and then follows $n-27$
from $47$ through $79$. Their interleaving explains the clusters.

\paragraph{Well-definedness.}
Suppose the prefix through $n-1$ is defined with $1\le b(k)\le k$.
For $1\le x\le n-1$ we have $1\le n-b(x)\le n-1$. Thus $T_n$ maps
the known index set into itself, and any positive finite iterate exists
uniquely. Strong induction proves well-definedness and
$1\le b(n)\le n-1$ for $n\ge2$. In fact, for $n\ge3$ the map takes
values in $[2,n-1]$: at $x=1$ its value is $n-1$, while at $x\ge2$ the
stronger bound $b(x)\le x-1$ gives $n-b(x)\ge n-x+1\ge2$.

\section{The candidate function and short orbits}
Define $F(1)=1$ and define $F(n)$ by the right side of
\eqref{eq:formula} for $n\ge2$. Its scale is unique because the respective
half-open intervals partition the eligible even and odd indices. Directly,
\begin{equation}\label{eq:bounds}
1\le F(n)\le n-1,\qquad F(n)\equiv n+1\pmod2\quad(n\ge2).
\end{equation}
The following four rules hold on closed intervals; the formulas agree at
every shared endpoint. In each row $t$ is a power of $3$.
\begin{center}
\begin{tabular}{lll}
\toprule
Parity of $x$ & Interval & $F(x)$\\
\midrule
even & $2t\le x\le4t$ & $x-t$\\
even & $4t\le x\le6t$ & $3t$\\
odd & $3t\le x\le5t$ & $2t$\\
odd & $5t\le x\le9t$ & $x-3t$\\
\bottomrule
\end{tabular}
\end{center}
At $x=6t$, for instance, the next even scale $3t$ gives the same value
$3t$; at $x=9t$ the next odd scale gives $6t$.

Fix $n$ and put $U_n(x)=n-F(x)$ on $1\le x\le n-1$. By
\eqref{eq:bounds}, this interval is invariant. Write $x_0=n-1$ and
$x_{j+1}=U_n(x_j)$.

\begin{lemma}\label{lem:orbit}
Let $s\ge3$ be a power of $3$ in the interval associated with $n$.
The orbit is given by the tables below. In particular, $x_6=x_4$;
if $n$ is even then $x_4=F(n)$, and if $n$ is odd then $x_5=F(n)$.
\end{lemma}

In the first table set
\[
A=\min(n-2s/3,5s/3+1),\qquad B=\min(n-2s/3,2s).
\]
Rows are closed for convenience, and are intersected with $2s\le n<6s$
or $3s\le n<9s$, respectively. Adjacent rows agree at shared endpoints.

\begin{center}
\textbf{Even $n$}\par\smallskip
\setlength{\tabcolsep}{3pt}
\resizebox{\textwidth}{!}{%
\begin{tabular}{c c c c c c c c}
\toprule
& Interval for $n$ & $x_1$ & $x_2$ & $x_3$ & $x_4$ & $x_5$ & $x_6$\\
\midrule
E1 & $[2s,3s+1]$ & $s+1$ & $n-2s/3-1$ & $A$ & $n-s$ & $B$ & $n-s$\\
E2 & $[3s+1,10s/3]$ & $n-2s$ & $7s/3$ & $n-4s/3$ & $n-s$ & $2s$ & $n-s$\\
E3 & $[10s/3,4s]$ & $n-2s$ & $n-s$ & $2s$ & $n-s$ & $2s$ & $n-s$\\
E4 & $[4s,5s+1]$ & $n-2s$ & $3s$ & $n-2s$ & $3s$ & $n-2s$ & $3s$\\
E5 & $[5s+1,6s]$ & $3s+1$ & $n-2s-1$ & $n-2s$ & $3s$ & $n-2s$ & $3s$\\
\bottomrule
\end{tabular}}
\end{center}

\begin{center}
\textbf{Odd $n$}\par\smallskip
\setlength{\tabcolsep}{3pt}
\resizebox{\textwidth}{!}{%
\begin{tabular}{c c c c c c c c}
\toprule
& Interval for $n$ & $x_1$ & $x_2$ & $x_3$ & $x_4$ & $x_5$ & $x_6$\\
\midrule
O1 & $[3s,4s+1]$ & $s+1$ & $n-2s/3-1$ & $5s/3+1$ & $n-s$ & $2s$ & $n-s$\\
O2 & $[4s+1,13s/3]$ & $n-3s$ & $10s/3$ & $n-7s/3$ & $n-s$ & $2s$ & $n-s$\\
O3 & $[13s/3,5s]$ & $n-3s$ & $n-s$ & $2s$ & $n-s$ & $2s$ & $n-s$\\
O4 & $[5s,6s+1]$ & $n-3s$ & $4s$ & $n-3s$ & $4s$ & $n-3s$ & $4s$\\
O5 & $[6s+1,7s]$ & $3s+1$ & $n-2s-1$ & $n-3s$ & $4s$ & $n-3s$ & $4s$\\
O6 & $[7s,8s+1]$ & $3s+1$ & $n-2s-1$ & $n-3s$ & $n-3s$ & $n-3s$ & $n-3s$\\
O7 & $[8s+1,9s]$ & $3s+1$ & $n-2s-1$ & $5s+1$ & $n-3s$ & $n-3s$ & $n-3s$\\
\bottomrule
\end{tabular}}
\end{center}

\begin{proof}
Use the four branch rules with $t=s/3,s,3s$, all powers of $3$. The
parities alternate for even $n$, starting with odd $x_0$; for odd $n$
all indices $x_j$ are even. All listed arguments remain in $[1,n-1]$.

We detail E1, which involves the two minima. Here
$5s/3\le2s-1\le n-1\le3s$, so the odd linear branch at scale $s/3$
gives $F(n-1)=n-s-1$ and $x_1=s+1$. Since
$2s/3\le s+1\le4s/3$, we get $F(x_1)=2s/3+1$ and
$x_2=n-2s/3-1$. This odd index lies in $[s,3s]$, and hence
\[
F(x_2)=\max(2s/3,n-5s/3-1),\qquad x_3=A.
\]
Both arguments defining $A$ are at least $4s/3$, and
$A\le5s/3+1\le2s$. The even plateau at scale $s/3$ therefore gives
$F(x_3)=s$ and $x_4=n-s$. This odd index lies in $[s,3s]$, so
$F(x_4)=\max(2s/3,n-2s)$ and $x_5=B$. Finally,
$4s/3\le B\le2s$ gives $F(x_5)=s$ and $x_6=n-s$.

The remaining rows follow by the same substitutions. To illustrate the
change of scale in O7, $x_1=3s+1$ gives $F(x_1)=2s+1$, so
$x_2=n-2s-1\in[6s,7s-1]$. The even linear branch at scale $3s$
gives $F(x_2)=n-5s-1$ and $x_3=5s+1$. The even plateau at scale $s$
then gives $x_4=n-3s$; this lies in $[5s+1,6s]$ and is fixed by $U_n$.
Each of the other rows uses the displayed branch intervals directly.
\end{proof}

\section{Proof of the formula}
We prove $b(n)=F(n)$ by strong induction. Direct evaluation through $n=8$
gives $(1,1,2,3,2,3,4,5)$, as required. For $n\ge9$, assume equality
at every smaller index. Then $T_n=U_n$ throughout $[1,n-1]$. Since this
interval is invariant, their entire starting orbits agree.

Let $d=b(n-1)=F(n-1)$. Equation~\eqref{eq:bounds} gives
$d\equiv n\pmod2$. For even $n\ge10$, the table gives $d\ge4$:
in E1, $d=n-s-1\ge n/2-1\ge4$; in E2--E4, $d=2s\ge6$;
in E5, $d=n-3s-1\ge2s\ge6$. For odd $n$, it gives $d\ge5$:
in O1, $d=n-s-1\ge2s-1\ge5$; in O2--O4, $d=3s\ge9$;
in O5--O7, $d=n-3s-1\ge3s\ge9$.

Since $x_6=x_4$, determinism implies $x_{j+2}=x_j$ for every $j\ge4$.
For even $n$, $d$ is even and at least four, so
$b(n)=x_d=x_4=F(n)$. For odd $n$, $d$ is odd and at least five, so
$b(n)=x_d=x_5=F(n)$. This closes the induction.

The same lemma proves~\eqref{eq:orbit} for $n\ge9$. For
$n=2,3,4,5,6,7,8$, direct calculation gives $x_4=x_6$ equal to
$1,2,3,4,3,4,5$, respectively. The bound four is attained at $n=24$:
\[
23\longmapsto10\longmapsto17\longmapsto16\longmapsto15
\longmapsto18\longmapsto15\longmapsto18\longmapsto\cdots.
\]
The first six displayed values are distinct. This finishes the proof of
Theorem~\ref{thm:main}.\hfill$\square$

\section{Scaling, fluctuations and fixed points}
\begin{corollary}
For every $n\ge2$, $b(3n)=3b(n)$. Moreover,
\[
\liminf_{n\to\infty}\frac{b(n)}n=\frac25,\qquad
\limsup_{n\to\infty}\frac{b(n)}n=\frac34.
\]
\end{corollary}
\begin{proof}
Multiplying $n$ by $3$ preserves parity, sends its scale $s$ to $3s$,
and multiplies either expression in~\eqref{eq:formula} by $3$.
On even indices, $b(n)/n\in[1/2,3/4]$; on odd indices it lies in
$[2/5,2/3]$. The values $n=5\cdot3^k$ attain $2/5$, and
$n=4\cdot3^k$ attain $3/4$, proving the two limits.
\end{proof}

For even $n\ge4$, the eventual period is exactly two, since the two
phases have opposite parity; at $n=2$ the orbit is fixed. For odd $n$,
the tables give
\[
x_4=\begin{cases}
n-s,&3s\le n\le5s,\\
4s,&5s\le n\le7s,\\
n-3s,&7s\le n<9s,
\end{cases}
\qquad x_5=\max(2s,n-3s).
\]
Consequently the prescribed orbit is eventually fixed exactly when
$n=3s$ or $7s\le n<9s$. At the other odd indices its period is two.
The small scale $s=1$ agrees by direct calculation at $n=3,5,7$.

\section{Computational support}
The accompanying package provides two complementary checks. The program
\texttt{explore.py} generates the original recurrence through $10^6$ terms
without using the closed formula, compares every term with it, and separately
performs literal nesting through $5000$ terms. All comparisons agree.
The program \texttt{check\_proof.py} checks the twelve orbit rows using exact
rational Fourier--Motzkin elimination, including the strict inequalities
needed to negate an implication. It verifies all 72 transitions over their
full parameter regions, together with interval coverage, index bounds,
parity, selected outputs, depth bounds and the small cases.

These computations supplement the written induction proof. This result has
not yet been formalized in Lean. AI assistance was used for exploration,
proof development, code generation and writing; the package records the
arithmetic checks performed on the resulting argument.

\paragraph{Acknowledgements.}
We thank John M. Campbell for sharing the recurrence and Beno\^it Cloitre
for the discussion of stabilization in nested iterations.

\end{document}
